\section{Strength and Weakness}

With the model complete, we set out the properties that make it useful to an archaeologist and then the limitations that bound how far its answers should be pushed.

\textbf{Integration of contributing factors.} The model gathers mechanical wear, the habits of the people who used the stairs, and the properties of the construction material into a single framework, so that wear is not assigned to any one cause in isolation. Because the damage from these sources accumulates within one scheme, the model presents how wear develops over the life of a stairwell more coherently than a chain of separate estimates could.

\textbf{A statistical account of footfalls.} Individual footfalls enter the model as random events, not as a fixed and endlessly repeated loading. A normal distribution, extended to the multivariate case, describes where each footprint falls along the two in-plane directions of the tread. Treating the footfalls in this way is what lets the model reflect how people actually walked, separating single-file passage from side-by-side passage and giving the two directions of travel the different weights they deserve.

\textbf{Measurements that leave the artefact intact.} Every quantity the model requires, among them the depth of the wear and the initial hardness of the stone, can be obtained without damaging the structure. The method therefore sits comfortably with the archaeological duty to protect the objects under study.

\textbf{Reliance on simplifying assumptions.} The formulation assumes that the factors driving wear stay broadly constant across long spans of time. Real buildings are not used so steadily, since traffic may surge during a season of pilgrimage and then fall silent for decades, and variation of that kind lies outside the model.

\textbf{A coarse treatment of environmental change.} A single uniform rate constant $p$ stands for the pace of corrosion and wear under every combination of humidity and temperature. The microclimate within one building and the cycling of the seasons in fact produce rates that vary in both space and time, and that finer layer has been simplified away.

Taken as a whole, the model is most credible for stairs whose traffic and surroundings stayed comparatively steady, and least credible where use or climate shifted sharply over the life of a structure.

\section{Conclusion}

To put better tools in the hands of archaeologists, we built the Stair Wear Model, which begins from a measured wear matrix of the stairs under study and answers the archaeologist's questions from it. The model has two parts. The Wear Volume Model (WVM) accounts for the amount of material lost from each step and expresses the relationships among the factors responsible for that loss. The Wear Distribution Model (WDM) accounts for the distribution of the wear across the tread, and from that distribution we examine how people once moved across these steps. Beyond these two components, we also provided consistency and reliability indicators for the wear distribution and for the age of the stairs. We set out further methods for judging the source of the stone and for deciding whether a set of stairs has been repaired or renovated. Simulation was used extensively in arriving at all of this.

Wear on a staircase cannot be avoided, and it is one of the few records a building keeps of the people who passed through it. Those who came before us left the mark of their footsteps in the stone, and that mark can be read.

The direction we would pursue next is the introduction of more detailed environmental variables, so that environmental conditions interfere less with the analysis of the stone. Our hope is that this model will help archaeologists recover the information they seek from the stairs they study.
